\documentclass[11pt,a4paper]{article}
\usepackage[T1]{fontenc}
\usepackage[utf8]{inputenc}
\usepackage{lmodern,amsmath,amssymb,booktabs,longtable,array,xcolor,microtype}
\usepackage[a4paper,left=2.3cm,right=2.3cm,top=2.5cm,bottom=2.5cm]{geometry}
\usepackage{upgreek}
\usepackage[colorlinks=true]{hyperref}
\definecolor{ojacblue}{rgb}{0,0,0.8}
\hypersetup{linkcolor=ojacblue,citecolor=ojacblue,urlcolor=ojacblue}
\renewcommand{\baselinestretch}{1.08}
\setlength{\parindent}{0pt}
\setlength{\parskip}{4pt}
\setlength{\emergencystretch}{3em}
\newcommand{\R}{https://github.com/ralfanop/Certifications_Eulerian_polar-curve_computations}
\newcommand{\PR}[1]{\href{\R/pull/#1}{PR~#1}}
\newcommand{\CI}[1]{\href{\R/actions/runs/#1}{CI~\texttt{#1}}}
\newcommand{\eq}[1]{(#1)}
\newcommand{\rE}{r_E}
\begin{document}

\begin{center}
{\Large\bfseries Manuscript 14: Lean~4 and Arb certification record}\\[6pt]
{\large Magnetic properties of the Eulerian profile at the special point $i/\sqrt2$, version 14.119}\\[6pt]
Roberto Alfano, 9~November~2026
\end{center}

\section*{1. Scope.}
This record states, for each numbered result of manuscript~14, version~14.119, and for the computed values added to its text on 8 and 9~November~2026, what is kernel-checked in Lean~4, which scalar values are enclosed by Arb at 2000 bits, which steps are classical inputs and which formalisation tasks remain open. Equation, section and result numbers are those of version~14.119. Up to \eq{175} and up to Theorem~10.1 they are those of versions~14.105 and 14.118: up to \eq{141} they are those of version~14.91, and from \eq{142} on they exceed those of version~14.91 by one, because version~14.95 added the line \eq{142} with fifty decimals of the root $x_0$ of \eq{139}. Version~14.119 adds Proposition~10.2 with the displays \eq{176}--\eq{177}, so that Remark~10.2 of version~14.118 is Remark~10.3. All certificates are in the main branch of the repository \url{\R} at commit \href{\R/tree/bf418e4e708bdf6ec3d1498df57f4061a60415ae}{\texttt{bf418e4e}}, and were audited from 5 to 9~November~2026 (those of versions~14.95--14.105 in \PR{63}, those of version~14.118 in \PR{64}, those of version~14.119 in \PR{69}); each pull request (PR) below links its workflow runs. The Lean certificates use Lean~\texttt{v4.35.0-rc3} with the pinned Mathlib dependency recorded in their workflows, compiled from source; every axiom audit lists only \texttt{propext}, \texttt{Classical.choice} and \texttt{Quot.sound}, with no \texttt{sorryAx} and no \texttt{Lean.ofReduceBool}. The numerical certificates use \texttt{python-flint~0.9.0}, exact rational arithmetic and Arb balls at 2000 bits; they certify the stated scalar quantities and do not prove infinite-dimensional operator identities.

\section*{2. Overview.}
In Table~1, ``yes'' means that the stated content is certified, apart from the classical inputs listed in Section~4, and ``---'' that it is certified numerically only.

\vspace{10pt}
{\small
\begin{longtable}{@{}>{\raggedright\arraybackslash}p{0.38\textwidth}>{\raggedright\arraybackslash}p{0.11\textwidth}>{\raggedright\arraybackslash}p{0.13\textwidth}>{\raggedright\arraybackslash}p{0.30\textwidth}@{}}
\caption*{\normalsize Table 1. Certification status of the numbered results.}\\[2pt]
\toprule
Result & Lean 4 & Arb & Pull requests\\
\midrule
\endfirsthead
\toprule
Result & Lean 4 & Arb & Pull requests\\
\midrule
\endhead
\bottomrule
\endfoot
Proposition 2.1 (special-value chain) & --- & yes & \PR{38}\\
Section 2.1 (discriminant of $F_n$) & yes & yes & \PR{62}\\
Corollary 2.2 (rational bounds) & --- & yes & \PR{38}\\
Lemma 2.3 (uniqueness of $z_0$) & yes & yes & \PR{44}\\
Proposition 3.1 (Sonine) & yes & yes & \PR{31}, \PR{38}\\
Proposition 3.3 (modified Helmholtz) & yes & yes & \PR{38}, branch \texttt{cor36}\\
Proposition 4.1 (Langevin) & yes & yes & \PR{38}, \PR{46}\\
Theorem 4.2 (noninteracting magnet) & yes & yes & \PR{40}, \PR{48}, \PR{61}\\
Section 4.5 and its Table 1 (superparamagnetic fields) & --- & yes & \PR{64}\\
Figure 1 (caption) and the paragraph of Figure 2 & yes & yes & \PR{64}\\
Figures 1--3 (plotted data) & --- & yes & \PR{64}\\
Displays (54)--(55) (Riccati, Painlev\'e~III) & --- & yes & \PR{61}\\
Display (60) ($n$-vector) & $\lambda=0$ & $\lambda=0,\pm\tfrac12$ & \PR{40}, \PR{47}\\
Definition 5.1 (Gram operator) & yes & --- & \PR{33}\\
Theorem 5.2 (Gram commutation) & yes & --- & \PR{37}\\
Corollary 5.3 (Eulerian specialisation) & yes & yes & \PR{33}, \PR{38}\\
Proposition 5.4 (spheroidal basis) & yes & --- & \PR{41}, \PR{42}, \PR{43}\\
Proposition 6.2 (two determinants) & yes & yes & \PR{44}, branch \texttt{dirichlet-finite}\\
Proposition 6.3 (resolvent traces) & yes & yes & \PR{44}\\
Corollary 6.4 (Hilbert--Schmidt trace) & yes & yes & \PR{44}\\
Proposition 6.5 (zeta expansion) & --- & yes & \PR{39}\\
Theorem 7.1 (modular duplication) & yes & yes & \PR{44}\\
Theorem 8.1 (quadratic covering) & yes & yes & \PR{45}\\
Proposition 8.2 (Montgomery--Taylor) & yes & yes & \PR{45}, \PR{61}\\
Proposition 8.4 (Dirichlet $L$-series, self-dual coupling) & yes & yes & \PR{45}, \PR{63}\\
Proposition 8.6 (paths and Vi\`ete) & yes & yes & \PR{45}, \PR{61}, \PR{62}, \PR{63}\\
Proposition 8.7 (classical \textit{PT}-symmetric oscillators) & yes & yes & \PR{61}, \PR{63}\\
Text after Proposition 8.7 (discriminants) & yes & yes & \PR{62}, \PR{63}\\
Proposition 8.8 (lattice walks) & yes & yes & \PR{45}\\
Lemma 9.2 (time reversal) & yes & yes & \PR{45}, \PR{62}\\
Theorem 9.3 (Bloch antipode) & yes & yes & \PR{45}, \PR{46}, \PR{62}\\
Corollary 9.4 (averages over pure states) & yes & yes & \PR{45}, \PR{46}\\
Proposition 9.5 (Brillouin splitting) & yes & yes & \PR{44}, \PR{45}\\
Text after Proposition 9.5 (electron spin) & --- & yes & \PR{62}\\
Proposition 9.7 (involutions) & yes & yes & \PR{45}\\
Theorem 10.1 (Eulerian--Parry) & yes & yes & \PR{30}\\
Proposition 10.2 (Parry index) & yes & yes & \PR{69}\\
\end{longtable}
}

\section*{3. Content of the certificates.}
\textbf{Sections 2--3.} The script \texttt{special\_point\_i\_sqrt2\_arb.py} encloses $\rE(z_0)=\sinh(a)/a$ and checks, in exact rational arithmetic, the bounds \eq{19}. The duplication identity \eq{15} is kernel-checked in the \texttt{Bridge} library. The Sonine certificates (\PR{31}; \CI{41853697588}) prove the integral moments $\mathrm c_n$ for every index, their factorial formula, the interchange of the integral with absolutely convergent even series, and the spherical integral identities of the series-defined $J_0$ and $I_0$ at every complex argument; the radial-series certificate (\CI{41852735785}) proves the intertwining for coefficient-defined radial operators under its summability hypothesis. The differential equation \eq{36} is kernel-checked in six declarations on the branch \texttt{manuscript14-cor36-lean-arb-20261104} (\CI{41824974185}).

\textbf{Section 4.} The values \eq{52}--\eq{53} are enclosed and proved correctly rounded (\CI{42207399251}). The module \texttt{BlochMeasureCertificate} (\PR{46}) proves $\int e^{x\omega_3}\,\mathrm d\mu=\sinh x/x$ with $\omega_3$ uniform on $(-1,1)$, which gives the normalisation \eq{38} through the coordinates of the first input of Section~2.2. The module \texttt{Theorem42Certificate} (\PR{48}; \CI{42356565741}; twenty-one audited statements) proves Theorem~4.2: for (i), $u\,\partial_u\log B(u)=[1+L(H/2)]/2$, its derivative $L'(H/2)/4$ in $H$, the mean of the quasi-power $B(u)^{n+1}$, and the mean and variance of any finite law tilted by $u^k$ as the first two derivatives of its logarithmic normalisation; for (ii), the probabilities $1/2$, $1/6$ and $1/4$ of a descent, of two adjacent descents and of two descents at distance at least two, for every $n$, hence \eq{46}, the mean $(n-1)/2$, \eq{47} and \eq{50}; for (iii), the factorisation into independent spins with the fields \eq{51}; and (iv). The script of \PR{40} (\CI{42239016609}; 145 checks) certifies \eq{46}--\eq{47} and all descent correlations for $n\leq9$, \eq{47} for $n\leq670$, \eq{50} with its coefficients, \eq{57}--\eq{59} at $a$ and $b$, the real zeros and \eq{51} for $n\leq18$, and \eq{60} at $\lambda=\pm1/2$. The module \texttt{PlanarSpinCertificate} (\PR{47}; \CI{42355540276}; eight audited statements) proves \eq{60} at $\lambda=0$: the arcsine law $\nu_0$ is the image of the uniform angle under $\cos$, its moments are $\binom{2m}{m}/4^m$ and $0$, and its transform is $\sum_k(x^2/4)^k/(k!)^2$, also as an average over the full circle; its script (\CI{42355345132}; 13 checks) compares the series with $I_0$ at ten points, among them $I_0(z_0)=J_0(1/\sqrt2)$, and with the integrals.

\textbf{Section 5.} \CI{41840633678} checks seven declarations on the complex-$L^2$ Gram integral, its exact parameter and measure, the nuclear decomposition and the basis-independent trace, and \CI{41841459971} thirteen kernel and graph-domain lemmas; the kernel balance, the core-extension criterion, the Gegenbauer spectrum and the trace formula \eq{64} are in the \texttt{GegenbauerFock} library. The module \texttt{GramCommutationTheorem52} (\PR{37}; \CI{42209958304}) proves Theorem~5.2 for every $c>0$, the inclusion $\mathcal S_{1/2,c}\mathcal H_{1/2}\subset\mathrm{Dom}(\mathcal A_c)$ and \eq{63}, for the graph closure of the polynomial Legendre operator perturbed by $M_{\xi^2}$. The module \texttt{GegenbauerSpectralRealisation} (\PR{42}; \CI{42260542889}) identifies this closure with \eq{8}--\eq{9}, proves it selfadjoint with compact resolvent, and proves that $\mathcal A_c+c^2+1$ has a compact inverse. The module \texttt{SpheroidalProposition54} (\PR{41}; \CI{42249877445}) proves Proposition~5.4: parity, simplicity within each parity class (from the moments rather than from \eq{74}), positivity of the $\mathcal S_{1/2,c}$-eigenvalues and a joint Hilbert basis with $\sum_n\mathrm g_n={}_1F_1(1/2;3/2;c)$. The module \texttt{SpheroidalDLMFIdentification} (\PR{43}; \CI{42272154201}) identifies the basis vectors with the spheroidal functions $\mathrm{Ps}_n^0(\cdot,-c^2)$, the definition of DLMF \S\S30.2--30.3 entering as an explicit hypothesis.

\textbf{Section 6.} The branch \texttt{manuscript14-dirichlet-finite-20261104} (\CI{41824926567}, \CI{41825497267}) checks fourteen declarations on the finite truncations of the products and traces, with thirteen Arb comparisons. The module \texttt{Section67ScalarCertificate} (\PR{44}; \CI{42283397911}; nineteen audited statements) proves the Euler products \eq{13}--\eq{14} as the products \eq{79}--\eq{80} over the spectra \eq{76} and \eq{78}, with \eq{81}, the spectral sums \eq{82}--\eq{84} and \eq{89}--\eq{90} for every $x>0$, the values \eq{85} and \eq{87}, and \eq{91}--\eq{92}. The script of \PR{39} (\CI{42224556937}; 92 checks) certifies Proposition~6.5: the Bernoulli numbers and the rational terms of \eq{97}, the values $\zeta(2m)$ with exact tail bounds, the termwise agreement of the two forms of \eq{94}, and the enclosure of $\log\rE(z_0)$.

\textbf{Section 7 and Lemma~2.3.} The same module proves Lemma~2.3 (uniqueness of $z_0$ in the upper half-plane, $S(z_0)=z_1=2z_0$, $S(z_1)=z_0$ and \eq{22}--\eq{23}) and Theorem~7.1: \eq{98}, the product \eq{99} over the spectrum of $A_2$ and its splitting into the products over the spectra of $A$ and $A_{\mathrm{DN}}$, \eq{100} in product form, \eq{101}--\eq{103} and the companion values. Its script (\CI{42282862670}; 45 checks) certifies \eq{21}--\eq{103} as ball identities, the products and traces with rigorous tails, and the rounding of the values of Remark~7.4.

\textbf{Sections 8--9.} The module \texttt{Section89Certificate} (\PR{45}; \CI{42299447127}; thirty-three audited statements) proves: in Theorem~8.1, the fibres $\{\tau,-\tau\}$ of $\mathcal Q$, the deck transformations, \eq{115} with $\rE$ the limit of the partial products of \eq{114}, $m(-H)=1-m(H)$, the lifts of \eq{116}, the doubling of winding numbers, the fibres over $\pm1/4$, \eq{118} and the uniqueness of $z_0$; Proposition~8.2 (\eq{121} from the Mittag-Leffler expansion of the cotangent, the logarithmic derivative, \eq{123}); in Proposition~8.4, \eq{126}, \eq{131}--\eq{135}, $4G_*^2+4G_*-1=0\Leftrightarrow\sinh2K_{\mathrm{sd}}=1$ and the arcsine moments; part~(i) of Proposition~8.6, \eq{136} with its zeros, the Perron eigenvalue, the eigenvector relation and \eq{137}; the two identities in the proof of Proposition~8.8; Lemma~9.2; $\omega_{\mathrm{Bl}}\circ\Theta=-\omega_{\mathrm{Bl}}$ and $|\omega_{\mathrm{Bl}}|=1$; \eq{161}, \eq{162} and the time reversal of the thermal state, the spectral forms \eq{159}--\eq{160} being \eq{80} and \eq{83} of \PR{44}; and parts (i), (iii) and (iv) of Proposition~9.7. The module \texttt{BlochMeasureCertificate} (\PR{46}; \CI{42307987424}; fifteen audited statements) proves that the image of $\mu$ under $\mathfrak z\mapsto(\omega_3,\arg\mathfrak z)$, the axial coordinate and the azimuth of $\omega_{\mathrm{Bl}}(\mathfrak z)$, is $(4\pi)^{-1}\mathrm d\omega_3\,\mathrm d\varphi$ on $(-1,1)\times(-\pi,\pi)$, that $\mu$ is a probability measure, \eq{156} for every real $x\neq0$ and \eq{157}, the mean $L(x)$ of $\omega_3$ under $e^{x\omega_3}\mathrm d\mu$, the invariance of the law of $\omega_3$ under $\mathfrak z\mapsto-1/\overline{\mathfrak z}$, and the bijection of $\mathbb C$ onto the sphere without its south pole. The script of \PR{45} (\CI{42299795864}; 115 checks) certifies the rounding of \eq{120}, \eq{125} and \eq{127}, the closed forms as ball identities, the traces of \eq{121} and \eq{123} with two-sided tail bounds, \eq{122} in exact arithmetic, \eq{128}--\eq{130}, the characteristic polynomials of the paths for $m\leq12$, \eq{138}, \eq{149} for $2n\leq24$, \eq{150}--\eq{151} at $q=1/8$, \eq{156}--\eq{157} as integrals, and \eq{158} against the sublevel sums for $J=1/2,1,3/2,7/2$.

\textbf{Additions of 7~November~2026 (evening).} The script \texttt{manuscript14\_\allowbreak evening\_\allowbreak additions\_\allowbreak arb.py} in \texttt{numerics/\allowbreak arb} (\PR{61}; \CI{42422614198}; 45 checks, exact rational and $\mathbb Q(\sqrt2)$ arithmetic and Arb at 2000 bits) certifies: the Riccati equation \eq{54} of Section~4.3 at six points, its general solution $L_\varrho$ for four values of $\varrho$, the saddle (eigenvalues $1$, $-2$, separatrix slope $1/3=L'(0)$) and the node (eigenvalues $1$, $2$) of the Riccati foliation, $\mathsf w(z)=1/z-\cot z$ with $\mathsf w(ix)=iL(x)$ and the third Painlev\'e equation \eq{55} at four complex points; in the proof of Theorem~4.2, Knuth's Bernoulli numbers ($B_1=1/2$), the equality of the kernel \eq{49} with the form of Borodin, Diaconis and Fulman, the generating function \eq{48}, the mean $(n-1)/2$ and the variance $(n+1)/12$ for $n\leq30$ and, with $\mathrm k(0)=-1/2$, the mean $-(n-1)/2$, the negative variance and the polynomials $3-u$ and $u^2-8u+13$, the Eulerian polynomials as $n!\det[I+(x-1)\mathcal K_{n-1}]$ for $n\leq7$, and Bernoulli's sums of powers with $B_1=\pm1/2$; the eight terms of \eq{122} and their zeta form; the middle term of \eq{138}; and Proposition~8.7: the roots of $32x^4-16x^2+1$, which is $32$ times \eq{139}, its irreducibility (Eisenstein at $2$ for $\mu_{\mathrm{hex}}^4-4\mu_{\mathrm{hex}}^2+2$ and the criterion for biquadratics) and primitivity, so that it is the minimal polynomial of $x_0$ over $\mathbb Z$, \eq{140}, the integral \eq{141}, Bender's polynomial \eq{143}--\eq{146} for seven rational values of $\gamma_{\mathit{PT}}^2\in[0,2)$, \eq{147} with $\cot(\pi/8)=1+\sqrt2$, \eq{148}, and the midradius of the square antiprism with unit edges. The module \texttt{QuarticResidueCertificate} (\PR{61}; \CI{42422614173}; fourteen audited statements) proves the algebraic identities of \eq{139}, \eq{140}, \eq{143}--\eq{148} for $x_0=\sqrt{2+\sqrt2}/\sqrt8$ and $x_1=1/(2\sqrt{2+\sqrt2})$, with \eq{139} in the integer form $32x^4-16x^2+1$; the integral \eq{141}, the irreducibility and the value of $\cot(\pi/8)$ are certified only in the script.

\textbf{Additions of 8~November~2026 (version 14.92).} The script \texttt{manuscript14\_\allowbreak v1492\_\allowbreak additions\_\allowbreak arb.py} in \texttt{numerics/\allowbreak arb} (\PR{62}; \CI{42478168014}; 30 checks, exact integer, rational and Gaussian-rational arithmetic and Arb at 2000 bits) certifies: in Section~2.1, that the discriminant of $F_n(t)=t^n-t^{n-1}-1$, in the normalisation $R(f,f')/a_n$ of Gelfand, Kapranov and Zelevinsky (Chapter~12, Eq.~(1.29)), is $(-1)^{n-1}[n^n+(n-1)^{n-1}]\neq0$, by Sylvester determinants for $n=2,\ldots,40$ and by their trinomial formula (1.38) for $n=2,\ldots,300$, the layout being checked against their formulas (1.33)--(1.35) on random polynomials; after Proposition~8.7, the discriminant $2^{23}$ of $32x^4-16x^2+1$, which is $32$ times \eq{139}, the discriminant $16(1-g^2)[4(1-g^2)-(2-\gamma^2)^2]^2$ of Bender's polynomial \eq{143} for 49 rational pairs, its value $s^{12}/128>0$ on \eq{145}--\eq{146} and its zeros where two frequencies coalesce or one vanishes; after Proposition~8.6, the dispersion relation of the lossless line of loops and $2\cos(\pi/8)=\mu_{\mathrm{hex}}$; after Proposition~9.5, the rounding of $a\,k_{\mathrm B}/\mu_{\mathrm B}\approx0.526$~T/K, of $3.80$~K at $2$~T, of $e^{1/\sqrt2}\approx2.028$ and of $3.88$~K, with $\mu_{\mathrm B}=9.274\times10^{-24}$~J/T and with the CODATA value; in Lemma~9.2, that $-i\sigma_2$ is the quaternion unit $e_2$ of the representation of Pandey, Kumar and Puri; in Theorem~9.3, that the Hopf projection $(2\overline{z_1}z_2,|z_1|^2-|z_2|^2)$ is the Bloch map \eq{152}. The module \texttt{V1492AdditionsCertificate} (\PR{62}; \CI{42478417105}; nine audited statements) proves the trinomial formula at the parameters of $F_n$ and its nonvanishing, the even case of the quartic formula, the discriminants of $32x^4-16x^2+1$ and of \eq{143}, the value on \eq{145}--\eq{146} and its positivity, and the factorisations at coalescing and vanishing frequencies; the formulas (1.35) and (1.38) themselves enter as classical inputs.

\textbf{Changes of versions 14.95--14.105 (8~November~2026).} The script \texttt{manuscript14\_\allowbreak v14104\_\allowbreak additions\_\allowbreak arb.py} in \texttt{numerics/\allowbreak arb} (\PR{63}; \CI{42487059795}; 34 checks, exact rational and $\mathbb Q(\sqrt2)$ arithmetic and Arb at 2000 bits) certifies: the fifty decimals of \eq{142}, exact by truncation (the fifty-first is $4$), with $x_0$ computed by four routes, $\mu_{\mathrm{hex}}/\sqrt8$, the integral of \eq{141} by rigorous Arb integration, $\sqrt{1+1/\sqrt2}/2$ and $1/(4\sin(\pi/8))$; the monic form \eq{139}, whose $32$-fold is $32x^4-16x^2+1$, its roots $\pm x_0$, $\pm x_1$, \eq{140}, \eq{147} and its discriminant $1/128=2^{23}/32^6$, by the Sylvester determinant, by the quartic formula (1.35) and by the squared differences of the roots; the forms $2(2-\gamma_{\mathit{PT}}^2)$ of \eq{145} and $1-[(2-\gamma_{\mathit{PT}}^2)/\sqrt8]^2$ of \eq{146}, \eq{144} as $\mathsf P_{\mathit{PT}}(s_{\mathit{PT}}x)=s_{\mathit{PT}}^4$ times \eq{139} for seven rational values of $\gamma_{\mathit{PT}}^2\in[0,2)$, the discriminant $s_{\mathit{PT}}^{12}/128$ of $\mathsf P_{\mathit{PT}}$ as $s_{\mathit{PT}}^{12}$ times that of \eq{139}, the frequencies $s_{\mathit{PT}}x_k$ and \eq{148}; \eq{126}--\eq{127}; the fraction form of \eq{137}, the first three factors $\cos(\pi/4)\cos(\pi/8)\cos(\pi/16)$ of Vi\`ete's product; and, for \eq{128}--\eq{129}, $\zeta(1/2)<0$, the value $L_2(1/2,\upchi_1)$ as $2^{-1/2}\zeta(1/2,1/2)$ through the Hurwitz zeta function and as Arb's Dirichlet $L$-function of the principal character modulo $2$, both equal to $(1-2^{-1/2})\zeta(1/2)$ as in Table~22 of Mathar (whose entries $s=2,3,4$ are reproduced), hence $1-L_2(1/2,\upchi_1)/\zeta(1/2)=1/\sqrt2$ and $1/2-L_2(1/2,\upchi_1)/\zeta(1/2)=(\sqrt2-1)/2$, the integral of \eq{128}, and the first thirty digits of $1+G_*=(1+\sqrt2)/2$ against OEIS A174968. The module \texttt{V14104AdditionsCertificate} (\PR{63}; \CI{42487059874}; eleven audited statements) proves that $32$ times \eq{139} is $32x^4-16x^2+1$, the discriminant $1/128$ of \eq{139} and its relation to $2^{23}$, that $x_0=\sqrt{1+1/\sqrt2}/2$ is a root of \eq{139}, the forms \eq{145}--\eq{146}, \eq{144} under \eq{145}--\eq{146}, the discriminant of $\mathsf P_{\mathit{PT}}$ on \eq{145}--\eq{146} as $s_{\mathit{PT}}^{12}$ times that of \eq{139}, \eq{148}, the fraction form of \eq{137}, and \eq{128}--\eq{129} from the identity $L_2(s,\upchi_1)=(1-2^{-s})\zeta(s)$ and $\zeta(1/2)\neq0$, taken as hypotheses.

\textbf{Changes of version 14.118 (8~November~2026).} The script \texttt{manuscript14\_\allowbreak v14118\_\allowbreak section45\_\allowbreak figures\_\allowbreak arb.py} in \texttt{numerics/\allowbreak arb} (\PR{64}; \CI{42505194706}; 132 checks, exact rational arithmetic and Arb at 2000 bits) certifies: in Section~4.5, $\mu_{\mathrm p}=2.85\,\mu_{\mathrm B}$ and $x\approx0.0065$ at 1~T for the oxygen molecule at 293~K, the inductions of about 54~T and 108~T at $H_0$ and $H_1$ and of about 246~T at the edge $|H|=5/2+1/\sqrt2$ of the window of Figure~1(b), those of about 5.3~mT and 10.5~mT for the cobalt particle at 300~K and below 0.075~T up to $|H|=10$, and $\ln(10^{11})\approx25$; the twenty entries of Table~1; in the paragraph of Figure~2, the centres $m(H_0)=0.44156$, $m(H_1)=0.38590$ and the widths $L'(a)/4=0.08129$, $L'(b)/4=0.07562$ of the histograms for every $n=10,\ldots,670$, the only inflection point of \eq{44} at $H=0$ ($L''<0$ on $(0,\infty)$, by an exact Bernoulli bound on $(0,1/2]$, Arb bisection on $[1/2,5]$ and an exact exponential bound on $[5,\infty)$, with $L''$ odd) and the convexity at $H_0$, with second derivative $-L''(a)/8>0$; in the caption of Figure~1, $m(H)=1+1/(e^H-1)-1/H$, $m(-H)=1-m(H)$, the approach to the saturation values as $1/|H|$, the window centred at $H_0$, and the frame, guides, markers, labels and tangent of the figure source; every plotted point of Figures~1--3: the 2001 and 1601 points of the two panels of Figure~1, the 4008 centres and widths of Figure~2 for $n=3,\ldots,670$, correctly rounded from the exact Eulerian numbers, with the top vertex $p_{3,2}(1)=2/3$, and the 536 remainders of Figure~3 within $10^{-12}$ with every floor decided, the Parry roots being isolated by exact sign changes, with $h_E\approx0.9961$ and the endpoint and the left limit at $\beta_{31}$; and the 52 decimals of $\operatorname{Tr}(2/(2A+I))$ in \eq{112}, Remark~7.4. The module \texttt{V14118Section45Certificate} (\PR{64}; \CI{42505444820}; nine audited statements) proves that $L$ is odd, $m(-H)=1-m(H)$, $\coth(H/2)=(e^H+1)/(e^H-1)$, $m(H)=1+1/(e^H-1)-1/H$, the bounds $1-1/H<m(H)<1$ for $H>0$ and $0<m(H)<-1/H$ for $H<0$, $H[1-m(H)]=1-H/(e^H-1)$, and the limits $H[1-m(H)]\to1$ as $H\to+\infty$ and $-H\,m(H)\to1$ as $H\to-\infty$. Version~14.118 corrects two sentences of Section~4.5 (the window of Figure~1(b) at about 246~T; the whole curve of Figure~1(a), up to $|H|=10$, below 0.075~T) and the sign of the second derivative at $H_0$ in the paragraph of Figure~2.

\textbf{Changes of version 14.119 (9~November~2026).} Version~14.119 adds Proposition~10.2, certified as described under Section~10 below, and reduces Question~11.3 to its second part. It also corrects sentences of Section~4.4, of the proof of Theorem~4.2 and of the Abstract that identified the discrete Eulerian law with a continuous density or attributed to the finite normalisations $\mathrm A_n(u)$ identities of the quasi-power approximation: the tilted histogram is, up to normalisation, the sampling at the integers of the tilted density of $n+1$ uniform variables; the ratio $[(1+u)/2]^{n+1}$ of the duplication is that of the normalisations $B(u)^{n+1}$, the exact ratio $\mathrm A_n(u^2)/\mathrm A_n(u)=u\prod_j(u^2-\alpha_j)/(u-\alpha_j)$ following from the factorisation of part~(iii), kernel-checked in \texttt{Theorem42Certificate}; and \eq{44} is written $\mathbb E[k/(n+1)]=m(H)$, $m(H)$ being the Langevin curve. None of these corrections changes a certified statement.

\textbf{Section 10.} Theorem~10.1 is kernel-checked in full (\CI{41812736721}): the adjacent thresholds, the greedy digits, the rational orbit-defect bounds, the endpoint jump, the rational series of $\beta_E=\sqrt2\sinh(1/\sqrt2)$ (module \texttt{EulerianSinhSeries}), its transcendence and the absence of eventual periodicity; the finite word at $\beta_{31}$ is checked separately. The Hermite--Lindemann instance is compiled from source adapted, under its Apache~2.0 licence and with attribution, from Yuyang Zhao's \href{https://github.com/leanprover-community/mathlib4/pull/28013}{Mathlib pull request \#28013}, not yet merged on 7~November~2026. Proposition~10.2 is kernel-checked in the module \texttt{V14119ParryIndexCertificate} (\PR{69}; \CI{42601024063}; ten audited statements): the strict monotonicity of $t^m(t-1)$ on $[1,\infty)$, by which $\beta<\beta_{m+1}$ exactly when $\beta^m(\beta-1)<1$; the equivalence of this inequality with $m<\log[1/(\beta-1)]/\log\beta$, with equality only at a root; the ceiling \eq{176}; the greedy remainders $\beta^j(\beta-1)$ and the word $1\,0^{n-1}\,1$; and the indices $31$ and $187$ for every $\beta$ in the rational brackets of the two proofs, with the margins $F_{187}<-1/200$ and $F_{188}>1/100$. The script \texttt{manuscript14\_\allowbreak v14119\_\allowbreak parry\_\allowbreak index\_\allowbreak arb.py} in \texttt{numerics/\allowbreak arb} (\PR{69}; \CI{42601023857}; 23 checks, exact rational arithmetic and Arb at 2000 bits) encloses the quotient of \eq{176} in $(30,31)$ at $x=1/\sqrt2$ and in $(186,187)$ at $x=1/\sqrt8$, checks the rational bracket of $\rE(z_0)$ and the bracket of $\beta_E$ of Theorem~10.1, the index \eq{176} against its definition and the greedy word $1\,0^{n-1}\,1$ for $x=k/100$, $k=1,\ldots,175$, every floor decided, and the coefficients $6$, $1/5$ and $-3/10$ of \eq{177} by exact series arithmetic in $\mathbb Q$, with the remainder of \eq{177} below $x^2\log(6/x^2)/100$ at $x=10^{-1},\ldots,10^{-6}$ and equal, within $x^4\log^2(6/x^2)$, to the exact next-order term $-x^2[(19/3150)\log(6/x^2)+27/2800]$.

\section*{4. Classical inputs.}
The following enter the certificates as explicit hypotheses or are not reproved: the inputs of Section~2.2 of the manuscript, among them the surface measure $\mathrm d\xi\,\mathrm d\varphi/(4\pi)$ of $\mathbb S^2$ in axial coordinates; the identification of the Fredholm determinant and trace of a positive trace-class operator with the product and sum over its spectrum; the quasi-power approximation $\mathrm A_n(u)=B(u)^{n+1}[1+O(|H/(H+2\pi i)|^n)]$ with $u=e^H$ (Flajolet--Sedgewick, Example~IX.37); the real, simple, negative zeros of $\mathrm A_n$ (Frobenius); the determinantal structure of the descent process (Borodin--Diaconis--Fulman, Theorem~5.1); the $\mathrm j$-values of Lemma~2.3; the covering-space statements of Theorem~8.1(i) and the exact sequence \eq{117}; the definition of the spheroidal functions in DLMF \S\S30.2--30.3; the discriminant formulas (1.29), (1.35) and (1.38) of Gelfand, Kapranov and Zelevinsky, Chapter~12; and the identity $L_2(s,\upchi_1)=(1-2^{-s})\zeta(s)$ of Table~22 of Mathar, with $\zeta(1/2)\neq0$, both checked in Arb at $s=1/2$. The results listed as open formalisation tasks in the earlier versions of this record are classical and enter as inputs as well: the behaviour at the origin of the radial operators $\mathsf H_2$, $\mathsf H_3$ on even entire functions and the uniqueness of the regular solution of the radial equation, from the theory of regular singular points (Propositions~3.1 and~3.3); Archimedes' theorem, the uniform law on $[-1,1]$ of the axial component of a uniform point of $\mathbb S^2$ (Proposition~4.1); the identification of ratios of zeta-regularised determinants of Sturm--Liouville operators with Fredholm determinants, of Gel'fand and Yaglom in the form of Kirsten and McKane (Proposition~6.2), and the expansion of $\log(\sinh x/x)$ in the even zeta values from Euler's product (Proposition~6.5, whose values are certified in Arb); the decomposition of $L^2(0,2)$ into the parts odd and even under the reflection $t\mapsto2-t$ (Theorem~7.1); Abel's theorem on power series at the boundary (the value \eq{128} of Proposition~8.4); the spectrum $2\cos(j\pi/(m+1))$ of the path of $m$ sites, with characteristic polynomial $U_m(y/2)$, and Euler's beta integral $\int_0^\infty y^{v-1}(1+y)^{-1}\,\mathrm dy=\pi/\sin(\pi v)$, $0<v<1$ (part (ii) of Proposition~8.6, \eq{138}, and the integral \eq{141} of Proposition~8.7); Eisenstein's criterion (the irreducibility in Proposition~8.7); the Neumann series of the resolvent of the adjacency operator of $\mathbb Z^2$ for $|q|<1/4$ and P\'olya's generating function of returns (Proposition~8.8); the stereographic homeomorphism of $\mathbb C\cup\{\infty\}$ onto $\mathbb S^2$ (Theorem~9.3 and part (ii) of Proposition~9.7); the invariance of the Fubini--Study measure under the antipodal map (Corollary~9.4); and Lazarevi\'c's inequality $(\sinh x/x)^3>\cosh x$ for $x\neq0$, which is the sign of $L''$ on $(0,\infty)$ and hence the single inflection point of Figure~1, certified also in Arb. The statements of the additions of version~14.92 that are quotations of the literature (the line of loops of Solymar and Shamonina, Dyson's classes, the Berry phase and the Hopf bundle) are not reproved. In Proposition~10.2 the remainder term of \eq{177} rests on the Taylor expansions of $\sinh$ and $\log$ with remainder; its coefficients are certified exactly and its size numerically.

\section*{5. Open formalisation tasks.}
None. The statements listed here in the earlier versions of this record are classical results and enter as the inputs of Section~4. For Proposition~2.1 and Corollary~2.2 this record lists only the numerical certificates, and the data of Figures~1--3 are certified as numbers, in Arb.

\section*{6. Numerical jobs.}
The workflow \texttt{manuscript14-arb-2000bit.yml}, introduced by \PR{38}, runs the fourteen numerical scripts of the manuscript at 2000 bits: 572 checks in all, of which 162 in the nine original scripts (63 of the register of 4~November~2026, 35 Bessel, 15 Gram and resolvent, 13 Dirichlet, 28 Sonine and modified Helmholtz, 8 for Proposition~4.1), 92 for Proposition~6.5, 145 for Section~4.6, 45 for Sections~6--7, 115 for Sections~8--9 and 13 for \eq{60} at $\lambda=0$. The exact rational truncation bounds are those of the original scripts; only the ball precision was raised. The same workflow runs, since \PR{61}, \PR{62} and \PR{63}, the certificates of the additions of 7~November~2026 and of versions~14.92 and 14.95--14.105 of 8~November~2026, with 45, 30 and 34 checks, and, since \PR{64}, the certificate of version~14.118, with 132 checks, and, since \PR{69}, that of version~14.119, with 23 checks.

\end{document}
